\section{The easy-gate layer: exact two-qubit synthesis (parallel track)}
\label{sec:easy_gate_layer}

This section documents the route the project formalized first, following an
earlier draft of the paper: decompose an $n$-qubit unitary into arbitrary
two-qubit gates plus one-qubit Cliffords, and then eliminate the arbitrary
two-qubit gates by an exact two-qubit synthesis theorem (``Lemma 11'').

Every node here is proved, but the reader should know where it sits:
Theorem~\ref{thm:clifford_t_universal} does \emph{not} depend on any of it.
The check is mechanical: replacing the Lean proof of Lemma~\ref{lem:lemma11} by
\texttt{sorry} and rebuilding leaves the main theorem \texttt{sorry}-free,
whereas Theorem~\ref{thm:clifford_rz_from_lemma1} below becomes
\texttt{sorry}-dependent.  The layer is kept, and kept compiling, because the
quantitative bounds of Section~\ref{sec:bounds} are stated against it.

\subsection{The two-qubit cosine--sine decomposition}

\begin{lemma}
  \label{lem:cs_square_block}
  \lean{TwoControl.CosineSine.squareBlockCSD_exists}
  \leanok
  Every $4\times 4$ unitary $V$ factors as
  \[
    V = \operatorname{blockDiag}(P_0,P_1)\cdot
        \operatorname{csCore}(c_0,c_1,s_0,s_1)\cdot
        \operatorname{blockDiag}(Q_0,Q_1)
  \]
  with $P_i, Q_i$ one-qubit unitaries and nonnegative reals satisfying
  $c_i^2 + s_i^2 = 1$, where $\operatorname{csCore}$ is the block matrix
  $\begin{pmatrix} C & -S \\ S & C\end{pmatrix}$ built from the diagonal
  matrices $C = \operatorname{diag}(c_0,c_1)$, $S = \operatorname{diag}(s_0,s_1)$.
\end{lemma}

\begin{proof}
  \leanok
  Split $V$ into $2\times2$ blocks $A,B,C,D$.  Unitarity gives
  $A^\dagger A + C^\dagger C = I$, so diagonalizing $A^\dagger A$ produces a
  common orthonormal basis in which $A$ and $C$ become the diagonal cosine and
  sine matrices, and the remaining blocks follow from the column relations.
\end{proof}

\begin{lemma}
  \label{lem:cs_block_core}
  \lean{TwoControl.CosineSine.csBlockCore_eq_conditionalRy}
  \leanok
  If $c_i = \cos(\theta_i/2)$ and $s_i = \sin(\theta_i/2)$ then
  $\operatorname{csCore}(c_0,c_1,s_0,s_1)$ is the conditional rotation
  $R_y(\theta_0)\oplus R_y(\theta_1)$ acting on the first qubit conditioned on
  the second.
\end{lemma}

\begin{proof}
  \leanok
  Compare blocks; the identity is the definition of $R_y$ on each diagonal
  entry.
\end{proof}

\begin{lemma}[Cosine--sine decomposition, two qubits]
  \label{lem:cs_exists_2q}
  \lean{TwoControl.CosineSine.cosinesine_exists}
  \leanok
  For every two-qubit unitary $V$ there are one-qubit unitaries
  $P_0,P_1,Q_0,Q_1$ and angles $\theta_0,\theta_1$ with
  \[
    V = \operatorname{blockDiag}(P_0,P_1)\cdot
        \operatorname{condRy}(\theta_0,\theta_1)\cdot
        \operatorname{blockDiag}(Q_0,Q_1).
  \]
\end{lemma}

\begin{proof}
  \leanok
  \uses{lem:cs_square_block, lem:cs_block_core}
  Apply Lemma~\ref{lem:cs_square_block}, convert each nonnegative pair
  $(c_i,s_i)$ on the unit circle into an angle $\theta_i$, and rewrite the core
  with Lemma~\ref{lem:cs_block_core}.
\end{proof}

\subsection{Exact two-qubit synthesis}

\begin{lemma}
  \label{lem:oneq_diag}
  \lean{TwoControl.unitary_diag2_exists}
  \leanok
  Every one-qubit unitary $U$ can be written $U = V \operatorname{diag}(u_0,u_1) V^\dagger$
  with $V$ unitary and $\lvert u_0\rvert = \lvert u_1\rvert = 1$.
\end{lemma}

\begin{proof}
  \leanok
  The spectral theorem for $2\times2$ unitaries, from the project's shared
  diagonalization helpers.
\end{proof}

\begin{lemma}
  \label{lem:demux_2q}
  \lean{TwoControl.Clifford.lemma4_demultiplex_two_qubit}
  \leanok
  For one-qubit unitaries $V_0, V_1$ there are one-qubit unitaries $P, Q$ and
  angles $\alpha_0,\alpha_1$ with
  \[
    \operatorname{blockDiag}(V_0,V_1)
      = (I\otimes Q)\cdot \operatorname{condRz}(\alpha_0,\alpha_1)
        \cdot (I\otimes P).
  \]
\end{lemma}

\begin{proof}
  \leanok
  \uses{lem:oneq_diag}
  Diagonalize $V_1V_0^\dagger$ and split the resulting phases symmetrically;
  this is the two-qubit case of the demultiplexing identity.
\end{proof}

\begin{proposition}
  \label{prop:crz_pair}
  \lean{TwoControl.Clifford.controlled_rz_pair_uses_standard_gates}
  \leanok
  For all $\alpha_0,\alpha_1$, the conditional rotation
  $\operatorname{condRz}(\alpha_0,\alpha_1)$ is \emph{exactly} a circuit over
  $\{C(X), H, T, R_z(\theta)\}$.
\end{proposition}

\begin{proof}
  \leanok
  Split $(\alpha_0,\alpha_1)$ into its mean and half-difference.  The mean is a
  local $R_z$; the half-difference is realized by the standard
  two-$C(X)$ construction for a controlled $R_z$.  No approximation is needed
  because $R_z(\theta)$ is in the gate set.
\end{proof}

\begin{proposition}
  \label{prop:oneq_htrz}
  \lean{TwoControl.Clifford.one_qubit_exact_h_t_rz}
  \leanok
  Every one-qubit unitary is, up to global phase, exactly a circuit over
  $\{H, T, R_z(\theta)\}$.
\end{proposition}

\begin{proof}
  \leanok
  \uses{lem:ry_via_rz, lem:oneq_diag, def:global_phase}
  Prepare the first column with an $R_zR_y$ circuit, diagonalize the remaining
  unitary by Lemma~\ref{lem:oneq_diag}, and absorb the diagonal factor into a
  global phase times a final $R_z$.  Lemma~\ref{lem:ry_via_rz} together with
  $S = T^2$ and $S^\dagger = T^6$ expresses the $R_y$ factor in the allowed
  primitives.
\end{proof}

\begin{lemma}[Lemma 11]
  \label{lem:lemma11}
  \lean{TwoControl.Clifford.lemma11_two_qubit_synthesis}
  \leanok
  Every two-qubit unitary is, up to global phase, exactly a circuit over
  $\{C(X), H, T, R_z(\theta)\}$.
\end{lemma}

\begin{proof}
  \leanok
  \uses{lem:cs_exists_2q, lem:demux_2q, prop:crz_pair, prop:oneq_htrz,
        lem:ry_via_rz, def:global_phase}
  Write $U = P R Q$ by Lemma~\ref{lem:cs_exists_2q}, with $P,Q$ block diagonal
  and $R$ a conditional $R_y$ pair.  Synthesize $P$ and $Q$ with
  Lemma~\ref{lem:demux_2q} followed by Proposition~\ref{prop:crz_pair} and
  Proposition~\ref{prop:oneq_htrz}.  Rewrite $R$ using
  Lemma~\ref{lem:ry_via_rz} into a conditional $R_z$ pair with local Clifford
  conjugators and apply Proposition~\ref{prop:crz_pair}.  The accumulated
  scalars compose into a single global phase.

  The theorem is stated up to global phase rather than exactly.  An exact
  version would additionally require realizing an arbitrary unit scalar as a
  circuit over the same gate set, which is a separate obligation and is not
  needed anywhere downstream.
\end{proof}

\begin{lemma}
  \label{lem:2q_clifford_rz}
  \lean{TwoControl.Clifford.Universal.two_qubit_gate_has_clifford_rz_circuit}
  \leanok
  Lemma~\ref{lem:lemma11} restated in the vocabulary of
  Section~\ref{sec:circuits}, so that the layers below can use the gate-set
  predicates.
\end{lemma}

\begin{proof}
  \leanok
  \uses{lem:lemma11}
  Direct restatement.
\end{proof}

\subsection{Decomposition into easy gates}

\begin{lemma}
  \label{lem:easy_2q}
  \lean{TwoControl.Clifford.Universal.two_qubit_unitary_is_easy_gate}
  \leanok
  Every two-qubit unitary is a one-gate circuit over $\operatorname{EasyGate} 2$.
\end{lemma}

\begin{proof}
  \leanok
  \uses{def:gate_easy, def:embedded_two}
  The identity placement embeds the gate into itself, and the
  arbitrary-two-qubit clause of $\operatorname{EasyGate}$ accepts it.
\end{proof}

\begin{lemma}[Old Lemma 1]
  \label{lem:lemma1_easy}
  \lean{TwoControl.Clifford.Universal.lemma1_decomposition_to_easy_gate_set}
  \leanok
  For $n \ge 2$, every $n$-qubit unitary is \emph{exactly} a circuit over
  $\operatorname{EasyGate} n$: arbitrary two-qubit gates together with embedded
  $H$, $S$, $S^\dagger$, $R_z(\theta)$.
\end{lemma}

\begin{proof}
  \leanok
  \uses{lem:cs_step, lem:cry_via_crz, lem:demux, lem:mottonen, lem:easy_2q,
        def:gate_easy}
  Recursively apply the cosine--sine step (Lemma~\ref{lem:cs_step}), convert
  each uniformly controlled $R_y$ factor into a uniformly controlled $R_z$
  factor (Lemma~\ref{lem:cry_via_crz}), demultiplex the block-diagonal side
  factors (Lemma~\ref{lem:demux}), and reduce controlled $R_z$ families with
  Lemma~\ref{lem:mottonen} until only one- and two-qubit gates remain.  The
  recursion stops at two qubits by Lemma~\ref{lem:easy_2q}.
\end{proof}

\begin{lemma}
  \label{lem:easy_factor}
  \lean{TwoControl.Clifford.Universal.easy_gate_factor_to_clifford_rz}
  \leanok
  Every single easy gate is synthesized up to global phase over
  $\{C(X), H, T, R_z(\theta)\}$.
\end{lemma}

\begin{proof}
  \leanok
  \uses{lem:2q_clifford_rz, prop:embedded_s, prop:embedded_sdag,
        def:gate_cliffordtrz}
  Five cases.  An arbitrary embedded two-qubit gate uses
  Lemma~\ref{lem:2q_clifford_rz}; $H$ and $R_z(\theta)$ are already allowed;
  $S$ and $S^\dagger$ expand as $T^2$ and $T^6$ by
  Propositions~\ref{prop:embedded_s} and~\ref{prop:embedded_sdag}.
\end{proof}

\begin{lemma}
  \label{lem:easy_circuit}
  \lean{TwoControl.Clifford.Universal.easy_circuit_to_clifford_rz}
  \leanok
  Every exact easy-gate circuit can be converted into a circuit over
  $\{C(X),H,T,R_z(\theta)\}$ synthesizing the same unitary up to global phase.
\end{lemma}

\begin{proof}
  \leanok
  \uses{lem:easy_factor}
  Replace each factor by Lemma~\ref{lem:easy_factor} and concatenate; the
  individual phases multiply into one.
\end{proof}

\begin{theorem}
  \label{thm:clifford_rz_from_lemma1}
  \lean{TwoControl.Clifford.Universal.clifford_rz_synthesis_from_lemma1}
  \leanok
  For $n \ge 2$, every $n$-qubit unitary is synthesized up to global phase over
  $\{C(X),H,T,R_z(\theta)\}$.
\end{theorem}

\begin{proof}
  \leanok
  \uses{lem:lemma1_easy, lem:easy_circuit}
  Compose Lemma~\ref{lem:lemma1_easy} with Lemma~\ref{lem:easy_circuit}.

  Compare Theorem~\ref{thm:clifford_rz_to_t_rz}, which reaches the same
  conclusion for all $n \ge 1$ through the paper's route without arbitrary
  two-qubit gates.  That is the version the main theorem uses.
\end{proof}
